\chapter{Eksperimen dan Evaluasi}
\label{chapter:eksperimen-evaluasi}

Bab ini mengevaluasi komponen akuisisi dan ekstraksi dokumen yang dirancang pada \autoref{chapter:analisis-rancangan}.
Evaluasi akuisisi mencakup \textit{throughput}, kegagalan, deduplikasi, dan fleksibilitas terhadap sumber baru, sedangkan evaluasi ekstraksi mencakup penentuan jalur ekstraksi per halaman dan pemulihan struktur hierarkis dokumen.
Urutan eksperimen ekstraksi mengikuti dependensi implementasinya.
Aturan \textit{routing} halaman ditetapkan terlebih dahulu, lalu digunakan untuk membentuk korpus evaluasi jalur \textit{native} dan jalur OCR pada \textit{legal parser}.

\section{Evaluasi \textit{Datasource} dan \textit{Scraper} Peraturan BPK}
\label{sec:evaluasi-datasource-scraper-bpk}

\subsection{Eksperimen \textit{Scraper} Portal Peraturan BPK}
\label{sec:eksperimen-scraper-bpk}

\subsubsection{Tujuan dan Pertanyaan Evaluasi}
\label{sec:tujuan-evaluasi-scraper-bpk}

Eksperimen ini mengevaluasi mekanisme akuisisi dokumen pada portal peraturan BPK.
Evaluasi difokuskan pada tiga ukuran keberhasilan Tujuan 1, yaitu \textit{throughput} ingesti, tingkat kegagalan, dan ketiadaan duplikasi penyimpanan.

Pertanyaan evaluasi yang dijawab oleh eksperimen ini adalah sebagai berikut.

\begin{enumerate}
  \item Bagaimana perubahan nilai \textit{concurrency} dan jeda permintaan memengaruhi durasi serta \textit{throughput} ingesti?
  \item Apakah konfigurasi dengan \textit{throughput} lebih tinggi menghasilkan lebih banyak \textit{failure} atau mengganggu kelengkapan enumerasi?
  \item Apa penyebab \textit{failure} yang terjadi, dan apakah penyebab tersebut berkaitan dengan konfigurasi atau dengan kondisi sumber?
  \item Apakah \textit{failure} pada satu entri dapat menghentikan pemrosesan entri lain atau keseluruhan \textit{run}?
  \item Apakah dokumen dengan isi identik hanya menghasilkan satu salinan \textit{payload} di \textit{storage} meskipun memiliki \textit{source identifier} yang berbeda?
\end{enumerate}

\subsubsection{Rancangan Eksperimen}
\label{sec:rancangan-eksperimen-scraper-bpk}

Portal peraturan BPK dipilih sebagai sumber eksperimen utama karena menyediakan dokumen peraturan perundang-undangan yang menjadi korpus domain hukum pada tugas akhir ini.
Google Drive tidak digunakan sebagai sumber eksperimen karena hanya disiapkan untuk evaluasi integrasi terhadap sumber kedua.
Portal BPK juga tidak mencantumkan batas laju permintaan publik yang terdokumentasi.
Keadaan tersebut membuat pengaruh \textit{concurrency} dan jeda permintaan perlu ditentukan melalui eksperimen, bukan melalui satu nilai batas yang telah ditetapkan oleh sumber.

Korpus eksperimen diperoleh dari hasil pencarian portal dengan \textit{tag} Pendidikan dan pendidikan dan pelatihan.
Korpus tersebut digunakan untuk keenam konfigurasi sehingga perbandingan antarkonfigurasi dilakukan pada himpunan entri yang sama.
\textit{Datasource} hanya menerima berkas PDF.
Enumerasi dilakukan sampai portal mengembalikan halaman kosong dengan kebijakan penghentian \mbox{\texttt{stop\_mode=none}}.

Setiap konfigurasi dijalankan satu kali pada \textit{datasource} baru.
Seluruh \textit{run} dijalankan secara berurutan pada lingkungan yang sama, yaitu Docker Compose dengan satu \textit{scraper worker}, RabbitMQ dengan \textit{prefetch} satu, PostgreSQL, dan MinIO.
Tidak terdapat \textit{profile} aktif maupun proses \textit{pipeline} hilir selama eksperimen.
Versi kode, citra Docker, kredensial, jaringan, dan mesin dipertahankan sama untuk seluruh konfigurasi.
Dengan demikian, perubahan hasil dapat dikaitkan dengan parameter konfigurasi yang diuji, sementara faktor lingkungan dijaga tetap.

Variabel konfigurasi terdiri atas \path{download_concurrency}, dua parameter jeda enumerasi, dan dua parameter jeda pengunduhan.
Parameter jeda enumerasi adalah \path{enumerate_delay_min_seconds} dan \path{enumerate_delay_max_seconds}.
Parameter jeda pengunduhan adalah \path{download_delay_min_seconds} dan \path{download_delay_max_seconds}.
Metrik yang dicatat meliputi durasi \textit{run}, jumlah entri, jumlah dokumen baru, jumlah \textit{skip}, jumlah \textit{failure}, \textit{throughput}, kategori \textit{failure}, jalur perubahan status, dan hasil deduplikasi.

\begin{table}[htb]
  \centering
  \scriptsize
  \caption{Hasil enam konfigurasi eksperimen \textit{scraper} BPK}
  \label{tab:hasil-scraper-bpk}
  \begin{tabularx}{\textwidth}{@{}>{\raggedright\arraybackslash}X >{\centering\arraybackslash}p{2.0cm} c r r r@{}}
    \toprule
    \textbf{Konfigurasi} & \shortstack{\textbf{Jeda enumerasi}\\\textbf{dan unduh}} & \shortstack{\textbf{\texttt{download\_}}\\\textbf{\texttt{concurrency}}} & \textbf{Durasi (jam)} & \textbf{Dokumen/jam} & \textbf{\textit{Failure}} \\
    \midrule
    \path{bpk-00-baseline} & 1--3 detik & 1 & 14,14 & 618,5 & 16 \\
    \path{bpk-01-concurrency-5} & 1--3 detik & 5 & 3,58 & 2\,444,2 & 15 \\
    \path{bpk-02-concurrency-10} & 1--3 detik & 10 & 4,63 & 1\,887,0 & 15 \\
    \path{bpk-03-both-fast} & 0,5--1,5 detik & 1 & 7,39 & 1\,184,0 & 15 \\
    \path{bpk-04-moderate-stress} & 0,5--1,5 detik & 5 & 2,18 & 4\,007,4 & 15 \\
    \path{bpk-05-high-stress} & 0,5--1,5 detik & 10 & 1,75 & 4\,990,5 & 15 \\
    \bottomrule
  \end{tabularx}
\end{table}

\subsubsection{Perbandingan \textit{Throughput} dan Durasi}
\label{sec:hasil-throughput-scraper-bpk}

Keenam \textit{run} memproses 8.822 entri yang sama dan masing-masing menghasilkan 8.745 dokumen baru.
Perbedaan antarkonfigurasi terutama terlihat pada waktu yang diperlukan untuk mencapai hasil tersebut.
Konfigurasi \path{bpk-00-baseline} membutuhkan waktu 14,14 jam dengan \textit{throughput} 618,5 dokumen baru per jam.
Sebaliknya, konfigurasi \path{bpk-05-high-stress} selesai dalam 1,75 jam dengan \textit{throughput} 4.990,5 dokumen baru per jam.

\begin{figure}[htb]
  \centering
  \includegraphics[width=\textwidth]{experiments/scraper-bpk-factor-comparison.png}
  \caption{Perbandingan durasi dan \textit{throughput} pada konfigurasi \textit{scraper} BPK}
  \label{fig:scraper-bpk-factor-comparison}
\end{figure}

Perbandingan tersebut menunjukkan bahwa pengaruh \textit{concurrency} bergantung pada rentang jeda yang digunakan.
Pada jeda 1--3 detik, \textit{concurrency} lima menghasilkan durasi 3,58 jam dan \textit{throughput} 2.444,2 dokumen baru per jam, sedangkan \textit{concurrency} sepuluh menghasilkan durasi 4,63 jam dan \textit{throughput} 1.887,0 dokumen baru per jam.
Pada jeda 0,5--1,5 detik, \textit{concurrency} sepuluh justru lebih cepat daripada \textit{concurrency} lima, yaitu 1,75 jam dibandingkan 2,18 jam.
Oleh karena itu, \textit{concurrency} yang lebih tinggi tidak dapat dianggap selalu menghasilkan \textit{throughput} yang lebih tinggi.

\subsubsection{Bentuk Penyelesaian \textit{Run}}
\label{sec:progress-scraper-bpk}

\begin{figure}[htb]
  \centering
  \includegraphics[width=\textwidth]{experiments/scraper-bpk-completion-progress.png}
  \caption{\textit{Progress} penyelesaian entri dan akumulasi \textit{byte} tersimpan pada keenam \textit{run} BPK}
  \label{fig:scraper-bpk-completion-progress}
\end{figure}

\autoref{fig:scraper-bpk-completion-progress} memperlihatkan perkembangan pekerjaan berdasarkan waktu yang telah berlalu sejak setiap \textit{run} dimulai.
Panel kiri menghitung persentase entri yang telah mencapai status terminal, termasuk entri berhasil, \textit{skip}, maupun \textit{failure}.
Panel kanan menghitung akumulasi ukuran dalam \textit{byte} dari dokumen baru yang berhasil disimpan.
Karena panel kanan hanya memasukkan dokumen yang tersimpan, entri yang berakhir sebagai \textit{skip} atau \textit{failure} tidak menambah nilainya.
Selain itu, satu entri dapat memiliki berkas yang jauh lebih besar daripada entri lain.
Akibatnya, penyelesaian sejumlah entri tidak selalu menghasilkan kenaikan jumlah \textit{byte} yang sebanding.

\subsubsection{Hasil \textit{Run} Representatif dan Kategori \textit{Failure}}
\label{sec:hasil-failure-skip-scraper-bpk}

Seluruh \textit{run} berstatus \texttt{completed} dan menyelesaikan enumerasi pada himpunan 8.822 entri yang sama.
Status \texttt{completed} pada tingkat \textit{run} berarti enumerasi telah selesai, bukan berarti setiap entri berhasil disimpan sebagai dokumen baru.
Karena hasil \textit{skip}, jalur status, dan kategori \textit{failure} pada konfigurasi \textit{non-baseline} berulang, satu \textit{run} digunakan untuk menjelaskan pola umum.
\textit{Run} representatif tersebut adalah \mbox{\texttt{bpk-01-concurrency-5}}, yaitu konfigurasi yang kemudian direkomendasikan sebagai titik operasi.

Pada \textit{run} representatif, 19 entri dilewati karena isi berkas identik dan 43 entri dilewati karena tipe berkas tidak didukung.
Ke-43 entri tersebut terdiri atas 24 berkas \texttt{.doc}, 16 berkas \texttt{.docx}, dua berkas \texttt{.rar}, dan satu berkas \texttt{.rtf}.
Setiap entri dengan tipe yang tidak didukung mengikuti jalur \path{pending -> downloading -> downloaded -> skipped}.
Hal ini menunjukkan bahwa kebijakan tipe berkas diterapkan setelah berkas diterima dan ekstensi aktualnya diketahui.

\begin{figure}[htb]
  \centering
  \includegraphics[width=\textwidth]{experiments/scraper-bpk-skip-outcomes.png}
  \caption{Hasil \textit{skip} pada \textit{run} representatif \textit{scraper} BPK}
  \label{fig:scraper-bpk-skip-outcomes}
\end{figure}

Pada \textit{run} representatif terdapat 15 \textit{failure}.
Enam \textit{failure} berupa HTTP 404 pada tautan unduh, sedangkan sembilan lainnya berupa tautan unduh yang tidak tersedia.
Kelima \textit{run} lain menunjukkan sumber dan kategori \textit{failure} yang sama.
Dengan demikian, kedua kategori tersebut lebih tepat dipahami sebagai kondisi sumber daripada akibat perbedaan \textit{concurrency} atau jeda.

Konfigurasi \textit{baseline} memiliki satu perbedaan, yaitu 42 \textit{unsupported-file skip} dan 16 \textit{failure}, sedangkan konfigurasi lain memiliki 43 \textit{unsupported-file skip} dan 15 \textit{failure}.
Perbedaan tersebut berasal dari sumber dengan \textit{source identifier} \href{https://peraturan.bpk.go.id/Details/120830/perbup-kab-purbalingga-no-9-tahun-2019}{\path{/Details/120830/perbup-kab-purbalingga-no-9-tahun-2019}}.
Pada konfigurasi \textit{baseline}, proses menghabiskan percobaan ulang sebelum menerima berkas \texttt{.rtf}.
Pada lima konfigurasi lain, berkas yang sama berhasil diterima lalu dilewati karena tipe berkasnya tidak didukung.
Tidak ditemukan HTTP 429 maupun HTTP 500 pada keenam \textit{run}.
Anomali tersebut karena itu kemungkinan besar berkaitan dengan gangguan jaringan pada sisi klien yang menyebabkan \textit{timeout} atau \textit{retry exhaustion}, bukan akibat \textit{rate limiting} dari \textit{server}.

\begin{figure}[htb]
  \centering
  \subfloat[Kategori \textit{failure}]{\includegraphics[width=0.47\textwidth]{experiments/scraper-bpk-error-categories.png}}
  \hfill
  \subfloat[Pengulangan sumber \textit{failure}]{\includegraphics[width=0.47\textwidth]{experiments/scraper-bpk-error-stability.png}}

  \vspace{0.7\baselineskip}
  \subfloat[Tipe berkas yang tidak didukung]{\includegraphics[width=0.47\textwidth]{experiments/scraper-bpk-unsupported-file-types.png}}
  \caption{Kategori \textit{failure}, kestabilan sumber, dan tipe berkas yang tidak didukung pada \textit{scraper} BPK}
  \label{fig:scraper-bpk-error-analysis}
\end{figure}

Tidak ditemukannya HTTP 429 pada seluruh \textit{run} menunjukkan bahwa portal peraturan BPK tidak memperlihatkan mekanisme \textit{rate limiting} terhadap pola permintaan yang digunakan selama eksperimen.
Kesimpulan ini bersifat terbatas pada kondisi eksperimen dan tidak menyatakan bahwa portal tersebut tidak memiliki pembatasan laju pada kondisi lain.
Kegagalan pada tingkat dokumen juga tidak menghentikan pemrosesan entri lain karena setiap entri diproses secara terisolasi.
Hal ini berbeda dari kegagalan enumerasi terminal, yang tetap dapat menghentikan keseluruhan \textit{run} sesuai rancangan pada \autoref{sec:solusi-pemisahan-akuisisi} dan \autoref{sec:solusi-pengendalian-run}.

\subsubsection{Deduplikasi dan Identitas Sumber}
\label{sec:analisis-deduplikasi-scraper-bpk}

Pada tahap awal akuisisi, \textit{source identifier} digunakan untuk melacak entri yang diberikan portal BPK.
Identitas tersebut digunakan untuk mengambil halaman rincian, menghubungkan metadata, dan merekam riwayat pemrosesan entri.
Setelah berkas PDF berhasil diterima, sistem menghitung \textit{hash} atas seluruh isi berkas.
Nilai \textit{hash} kemudian digunakan untuk menentukan apakah \textit{payload} identik sudah pernah disimpan.

Pada \textit{run} representatif ditemukan 19 pasang \textit{source identifier}.
Setiap pasangan memiliki nilai \mbox{\texttt{content\_sha256}} yang sama.
Keanggotaan seluruh klaster tersebut stabil pada keenam \textit{run}, sehingga jumlah klaster dan hasil deteksi duplikasi dapat dijelaskan melalui satu \textit{run} representatif.
Deduplikasi berdasarkan \textit{hash} memastikan bahwa isi PDF yang identik tidak disimpan sebagai dua salinan \textit{payload}.

\begin{figure}[htb]
  \centering
  \subfloat[Klaster sumber dengan isi identik]{\includegraphics[width=0.47\textwidth]{experiments/scraper-bpk-duplicate-clusters.png}}
  \hfill
  \subfloat[Perubahan \textit{winner} penyimpanan]{\includegraphics[width=0.47\textwidth]{experiments/scraper-bpk-duplicate-winners.png}}
  \caption{Deduplikasi berdasarkan kesamaan isi pada eksperimen \textit{scraper} BPK}
  \label{fig:scraper-bpk-deduplication}
\end{figure}

Aspek yang berubah antarkonfigurasi bukan keanggotaan klaster, melainkan \textit{winner}, yaitu \textit{source identifier} yang lebih dahulu mencapai tahap penyimpanan.
Sebelas dari 19 klaster mengalami perubahan \textit{winner} pada setidaknya satu konfigurasi.
Perubahan tersebut dapat terjadi karena urutan hasil enumerasi portal tidak dijamin stabil dan \textit{concurrency} mengubah urutan penyelesaian pengunduhan.
Dengan demikian, variasi \textit{winner} tidak menunjukkan bahwa deduplikasi gagal.
Variasi tersebut menunjukkan bahwa sistem perlu memisahkan identitas sumber dari identitas isi berkas.

\subsubsection{Kesimpulan Eksperimen}
\label{sec:kesimpulan-evaluasi-bpk}

Kesimpulan eksperimen ini disusun untuk menjawab lima pertanyaan evaluasi pada bagian sebelumnya.

\begin{enumerate}
  \item Perubahan \textit{concurrency} dan jeda permintaan memengaruhi durasi serta \textit{throughput} secara bersamaan.
  Konfigurasi \mbox{\texttt{bpk-05-high-stress}} menghasilkan \textit{throughput} tertinggi, yaitu 4.990,5 dokumen baru per jam, tetapi hasil perbandingan menunjukkan bahwa \textit{concurrency} lebih tinggi tidak selalu lebih cepat.
  Konfigurasi operasional yang dipilih adalah \path{download_concurrency}=5 dengan jeda 1--3 detik karena memberikan \textit{trade-off} yang lebih konservatif, dengan durasi 3,58 jam dan \textit{throughput} 2.444,2 dokumen baru per jam.

  \item \textit{Throughput} yang lebih tinggi tidak mengurangi kelengkapan enumerasi atau meningkatkan \textit{failure} pada keenam \textit{run} yang diamati.
  Seluruh konfigurasi memproses 8.822 entri, menghasilkan 8.745 dokumen baru, dan berstatus \texttt{completed}.
  Jumlah \textit{failure} juga serupa pada konfigurasi non-\textit{baseline}; perbedaan pada \textit{baseline} disebabkan oleh satu anomali sumber.

  \item \textit{Failure} yang berulang terutama berasal dari kondisi sumber, yaitu HTTP 404 dan tautan unduh yang tidak tersedia.
  Satu \textit{retry exhaustion} pada \textit{baseline} kemungkinan berkaitan dengan gangguan jaringan pada sisi klien karena tidak ditemukan HTTP 429 maupun HTTP 500.
  Tipe berkas yang tidak didukung tidak dihitung sebagai \textit{failure}, melainkan sebagai \textit{skip} sesuai kebijakan \textit{datasource}.

  \item \textit{Failure} pada satu entri tidak menghentikan pemrosesan entri lain.
  Setiap entri diproses secara terisolasi, sehingga seluruh \textit{run} tetap dapat menyelesaikan enumerasi meskipun sebagian entri berakhir sebagai \textit{failure} atau \textit{skip}.

  \item Deduplikasi mencegah penyimpanan ulang terhadap isi yang identik.
  Setelah pemeriksaan awal berdasarkan \textit{source identifier}, \textit{content\_sha256} menjadi sumber kebenaran akhir untuk identitas isi.
  Eksperimen menemukan 19 klaster duplikasi; pada setiap klaster hanya satu dokumen disimpan, sedangkan entri lain dengan hash yang sama dianggap sebagai duplikat dan berakhir sebagai \textit{skipped}.
\end{enumerate}

\subsection{Analisis Data Korpus BPK}
\label{sec:data-mining-korpus-bpk}

\subsubsection{Ruang Lingkup dan Korpus Analisis Data}
\label{sec:ruang-lingkup-data-mining-bpk}

Analisis data dilakukan pada \textit{metadata} dokumen di PostgreSQL untuk \textit{datasource} \texttt{bpk-peraturan-pendidikan}.
Korpus terdiri atas 8.745 dokumen dengan total ukuran penyimpanan sekitar 20,07 GB.
Median ukuran berkas adalah 463,74 KB, dan tahun peraturan yang tercatat berada pada rentang 1947--2026.

\begin{table}[htb]
  \centering
  \footnotesize
    \caption{Ringkasan korpus pada analisis \textit{metadata} BPK}
  \label{tab:ringkasan-korpus-bpk}
  \begin{tabularx}{0.88\textwidth}{@{}>{\raggedright\arraybackslash}X r@{}}
    \toprule
    \textbf{Karakteristik} & \textbf{Nilai} \\
    \midrule
    Jumlah dokumen & 8.745 \\
    Total \textit{storage} & 20,07 GB \\
    Median ukuran berkas & 463,74 KB \\
    Rentang tahun peraturan & 1947--2026 \\
    Dokumen dengan relasi & 2.839 \\
    Jumlah relasi & 4.364 \\
    \bottomrule
  \end{tabularx}
\end{table}

\subsubsection{Komposisi Korpus Peraturan}
\label{sec:komposisi-korpus-bpk}

Empat panel pada \autoref{fig:scraper-bpk-corpus-composition} memperlihatkan komposisi korpus berdasarkan tahun peraturan, bentuk hukum, status keberlakuan, dan subjek.
Keempat distribusi tersebut digunakan untuk membaca periode dokumen, jenis peraturan, keadaan keberlakuan, serta fokus topik korpus yang berhasil disimpan.

\begin{figure}[htb]
  \centering
  \subfloat[Distribusi tahun]{\includegraphics[width=0.47\textwidth]{experiments/scraper-bpk-corpus-year.png}}
  \hfill
  \subfloat[Bentuk hukum]{\includegraphics[width=0.47\textwidth]{experiments/scraper-bpk-corpus-legal-form.png}}

  \vspace{0.7\baselineskip}
  \subfloat[Status keberlakuan]{\includegraphics[width=0.47\textwidth]{experiments/scraper-bpk-corpus-status.png}}
  \hfill
  \subfloat[Subjek peraturan]{\includegraphics[width=0.47\textwidth]{experiments/scraper-bpk-corpus-subject.png}}
  \caption{Komposisi korpus peraturan BPK berdasarkan \textit{metadata}}
  \label{fig:scraper-bpk-corpus-composition}
\end{figure}

Panel distribusi tahun menunjukkan bahwa dokumen dalam korpus berasal dari rentang 1947--2026.
Jumlah tertinggi berada pada tahun 2019 dan 2020, masing-masing sebanyak 817 dokumen.
Dengan demikian, korpus mencakup dokumen historis, tetapi konsentrasinya lebih besar pada peraturan yang relatif baru.

Panel bentuk hukum menunjukkan dominasi empat bentuk peraturan, yaitu \texttt{PERBUP}, \texttt{PERWALI}, \texttt{PERDA}, dan \texttt{PERGUB}.
Keempat bentuk tersebut mencakup 80,4\% dari seluruh dokumen.
Komposisi ini menunjukkan bahwa korpus terutama berisi peraturan yang berkaitan dengan pemerintahan daerah dan administrasi regional.

Status \texttt{BERLAKU} mencakup 87,6\% dokumen.
Status \texttt{TIDAK\_BERLAKU} mencakup 12,4\% lainnya.
Dokumen yang tidak berlaku tetap menjadi bagian yang terukur dalam korpus, sehingga distribusi tersebut memperlihatkan bahwa korpus tidak hanya berisi peraturan yang masih berlaku.

Panel subjek menunjukkan bahwa \texttt{PENDIDIKAN} tercatat pada 95,1\% dokumen.
Dominasi tersebut sesuai dengan \textit{tag} pencarian yang digunakan saat akuisisi, tetapi keberadaan subjek lain menunjukkan bahwa korpus tetap mencakup topik administrasi regional yang lebih luas.

\subsubsection{Kelengkapan \textit{Metadata}}
\label{sec:kelengkapan-metadata-bpk}

\autoref{fig:scraper-bpk-metadata} memperlihatkan dua aspek metadata, yaitu kelengkapan pengisian \textit{field} serta distribusi nilai lokasi dan T.E.U.
Panel pertama menunjukkan seberapa banyak dokumen memiliki nilai pada setiap \textit{field}, sedangkan panel kedua menunjukkan nilai administratif yang paling sering muncul.

\begin{figure}[htb]
  \centering
  \subfloat[Kelengkapan \textit{metadata}]{\includegraphics[width=0.47\textwidth]{experiments/scraper-bpk-metadata-completeness.png}}
  \hfill
  \subfloat[Lokasi dan T.E.U.]{\includegraphics[width=0.47\textwidth]{experiments/scraper-bpk-location-teu.png}}
  \caption{Kelengkapan dan variasi \textit{metadata} korpus BPK}
  \label{fig:scraper-bpk-metadata}
\end{figure}

Pada panel kelengkapan, \texttt{subjek} tersedia pada 100\% dokumen, sedangkan \texttt{bidang} hanya tersedia pada 25\% dokumen.
Perbedaan ini menunjukkan bahwa kelengkapan metadata tidak seragam antar-\textit{field}.
Untuk pengelompokan berdasarkan topik, \texttt{subjek} menyediakan cakupan yang lebih konsisten daripada \texttt{bidang}.

Panel lokasi dan T.E.U. menunjukkan 552 nilai berbeda pada kedua metadata administratif tersebut.
Nilai yang beragam menunjukkan bahwa korpus mencakup banyak lokasi atau instansi penerbit dan tidak terpusat pada satu nilai administratif.
Pada saat yang sama, distribusi frekuensi yang tidak seragam menunjukkan bahwa sebagian nilai muncul jauh lebih sering daripada nilai lainnya.

\subsubsection{Distribusi Ukuran Berkas dan \textit{Storage}}
\label{sec:ukuran-berkas-storage-bpk}

\autoref{fig:scraper-bpk-file-size} membandingkan jumlah dokumen dan kontribusi \textit{storage} pada setiap kelompok ukuran berkas.
Panel pertama menunjukkan komposisi jumlah dokumen, sedangkan panel kedua menunjukkan jumlah ruang yang digunakan oleh kelompok ukuran yang sama.

\begin{figure}[htb]
  \centering
  \includegraphics[width=0.92\textwidth]{experiments/scraper-bpk-file-size.png}
  \caption{Distribusi ukuran berkas pada korpus BPK}
  \label{fig:scraper-bpk-file-size}
\end{figure}

Sebanyak 13,9\% dokumen memiliki ukuran lebih dari 5 MB, tetapi kelompok tersebut menyumbang 65,9\% dari total \textit{storage}.
Perbedaan antara proporsi dokumen dan proporsi penyimpanan menunjukkan bahwa beban \textit{storage} tidak berbanding lurus dengan jumlah dokumen.
Karena itu, distribusi ukuran berkas perlu diperhitungkan ketika memperkirakan kapasitas penyimpanan korpus.

\subsubsection{Relasi Antarperaturan}
\label{sec:relasi-antarperaturan-bpk}

\autoref{fig:scraper-bpk-relations} memperlihatkan struktur relasi antarperaturan melalui jenis relasi, cakupan target, jumlah relasi keluar per dokumen, dan dokumen dengan jumlah relasi masuk tertinggi.
Grafik tersebut digunakan untuk membedakan relasi yang tercatat pada metadata dari relasi yang target dokumennya dapat ditemukan di dalam korpus.

\begin{figure}[htb]
  \centering
  \subfloat[Ringkasan jenis relasi]{\includegraphics[width=0.47\textwidth]{experiments/scraper-bpk-relation-overview.png}}
  \hfill
  \subfloat[Target dengan jumlah relasi masuk tinggi]{\includegraphics[width=0.47\textwidth]{experiments/scraper-bpk-relation-targets.png}}
  \caption{Relasi antarperaturan pada \textit{metadata} korpus BPK}
  \label{fig:scraper-bpk-relations}
\end{figure}

Korpus memuat 4.364 relasi yang berasal dari 2.839 dokumen.
Jenis relasi yang paling banyak muncul adalah pencabutan, dicabut oleh, perubahan, dan diubah oleh.
Keempat jenis tersebut menunjukkan bahwa metadata korpus tidak hanya menyimpan dokumen secara terpisah, tetapi juga mencatat hubungan perubahan dan keberlakuan antarperaturan.

Pada panel cakupan target, 1.000 sitasi tidak memiliki target BPK, sedangkan 888 target relasi berada di luar korpus yang dianalisis.
Relasi dengan target yang tersedia di dalam korpus merupakan bagian yang dapat dihubungkan langsung, sementara dua kategori lainnya tetap tercatat sebagai relasi tanpa dokumen target lokal.

Distribusi jumlah relasi per dokumen menunjukkan bahwa sebagian besar dokumen memiliki sedikit atau tidak memiliki relasi, sedangkan hanya sebagian kecil dokumen memiliki banyak relasi.
Panel dengan skala logaritmik menegaskan ekor panjang tersebut dengan memperlihatkan dokumen yang memiliki jumlah relasi jauh lebih tinggi daripada mayoritas dokumen.

Panel target relasi menampilkan peraturan yang paling sering menjadi tujuan relasi masuk.
Konsentrasi relasi pada sejumlah target menunjukkan bahwa beberapa peraturan menjadi rujukan atau objek perubahan bagi banyak dokumen lain.
Target yang tidak terdapat di dalam korpus tetap dihitung berdasarkan identitas targetnya, sehingga hasil analisis tidak menyamakan ketiadaan dokumen lokal dengan ketiadaan relasi pada sumber.

\section{Evaluasi \textit{Datasource} dan \textit{Scraper} Google Drive}
\label{sec:evaluasi-datasource-scraper-google-drive}

\subsection{Tujuan dan Pertanyaan Evaluasi}
\label{sec:tujuan-evaluasi-google-drive}

Evaluasi ini menilai fleksibilitas mekanisme akuisisi terhadap sumber baru menggunakan \texttt{google-drive-sanity}.
Berbeda dari eksperimen BPK, evaluasi ini tidak digunakan untuk memilih konfigurasi atau membandingkan \textit{throughput}.
Fokusnya adalah memeriksa apakah sumber Google Drive dapat masuk ke alur akuisisi yang sama melalui \textit{plugin} dan konfigurasi khusus sumber.

Pertanyaan evaluasi yang dijawab adalah sebagai berikut.

\begin{enumerate}
  \item Apakah Google Drive dapat menggunakan kontrak \textit{datasource} dan \textit{scraper} yang sama dengan sumber BPK?
  \item Apakah hasil enumerasi dan pengunduhan dari \textit{plugin} Google Drive dapat diteruskan ke validasi tipe berkas, penyimpanan, deduplikasi, dan pencatatan \textit{lifecycle} pada komponen inti?
  \item Apakah penambahan sumber Google Drive dapat dilakukan tanpa mengubah orkestrasi inti?
\end{enumerate}

\subsection{Rancangan Evaluasi}
\label{sec:rancangan-evaluasi-google-drive}

Evaluasi menggunakan satu folder Google Drive yang berisi sebelas entri, yaitu sepuluh PDF dan satu berkas ZIP.
\textit{Plugin} Google Drive mengakses folder melalui Google Drive API menggunakan \textit{service account}, melakukan enumerasi dengan \textit{pagination} standar, dan tidak menggunakan batas penghentian tambahan.
ID berkas Google Drive digunakan sebagai \textit{source identifier}, sedangkan nama berkas, tipe MIME, waktu perubahan, ukuran, dan tautan tampilan diteruskan sebagai metadata entri.

\textit{Datasource} dikonfigurasi hanya untuk menerima berkas PDF.
Kebijakan tipe berkas diterapkan oleh komponen inti, sehingga \textit{plugin} tidak membuat aturan ekstensi kedua yang dapat berbeda dari konfigurasi \textit{datasource}.
Evaluasi menggunakan jeda enumerasi dan pengunduhan 1--3 detik dengan \textit{concurrency} lima.
Konfigurasi tersebut dijalankan satu kali pada \textit{datasource} baru.

Evaluasi ini tidak dirancang sebagai eksperimen pembandingan.
Jumlah entri, ukuran berkas, API, mekanisme \textit{pagination}, dan model \textit{retry} Google Drive berbeda dari portal BPK, sehingga hasil waktunya tidak dibandingkan dengan \textit{throughput} BPK.

\subsection{Hasil Pemrosesan}
\label{sec:hasil-evaluasi-google-drive}

\textit{Plugin} Google Drive memproses seluruh sebelas entri dalam 33,0 detik.
Sembilan berkas PDF berhasil dibuat sebagai dokumen baru.
Satu berkas PDF dilewati karena isinya sudah direpresentasikan oleh dokumen yang tersimpan, satu berkas \path{netgear.zip} dilewati karena tipe berkasnya tidak diterima, dan tidak terdapat \textit{failure}.
Sembilan entri PDF yang menjadi dokumen baru mengikuti jalur \path{pending -> downloading -> downloaded -> completed}.
Entri PDF duplikat mengikuti jalur \path{pending -> downloading -> downloaded -> skipped}, sedangkan berkas ZIP mengikuti jalur \path{pending -> skipped} sebelum pengunduhan.

\begin{table}[htb]
  \centering
  \footnotesize
  \caption{Hasil evaluasi integrasi \textit{scraper} Google Drive}
  \label{tab:hasil-evaluasi-google-drive}
  \begin{tabular}{@{}lr@{}}
    \toprule
    \textbf{Metrik} & \textbf{Nilai} \\
    \midrule
    Total entri & 11 \\
    Dokumen baru & 9 \\
    \textit{Skip} duplikat & 1 \\
    \textit{Skip} tipe berkas & 1 \\
    \textit{Failure} & 0 \\
    Durasi & 33,0 detik \\
    \bottomrule
  \end{tabular}
\end{table}

\begin{figure}[htb]
  \centering
  \includegraphics[width=0.78\textwidth]{experiments/scraper-google-drive-latest-run.png}
  \caption{Hasil evaluasi integrasi \textit{scraper} Google Drive}
  \label{fig:scraper-google-drive-evaluasi}
\end{figure}

\subsection{Analisis Fleksibilitas Sumber}
\label{sec:analisis-fleksibilitas-google-drive}

Hasil pemrosesan menunjukkan bahwa mekanisme enumerasi Google Drive yang menggunakan API dapat dipetakan ke kontrak \textit{scraper} yang sama dengan portal BPK.
Perbedaan cara memperoleh daftar berkas dan mengunduh isi dokumen ditangani oleh \textit{plugin} Google Drive, sedangkan komponen inti tetap menangani validasi tipe, penghitungan \textit{hash}, deduplikasi, penyimpanan, dan pencatatan status.

Jalur sembilan PDF yang berakhir sebagai dokumen baru menunjukkan bahwa keluaran \textit{plugin} dapat diteruskan sampai tahap penyimpanan tanpa perubahan pada orkestrasi.
Entri \path{netgear-xr500.pdf} memiliki \textit{source identifier} yang berbeda dari \path{netgear-xr500 (1).pdf}, tetapi nilai \path{content_sha256}-nya sama.
Karena itu, entri tersebut dianggap sebagai duplikat setelah pengunduhan dan berakhir sebagai \textit{skipped}.

Jalur \path{pending -> skipped} pada \path{netgear.zip} menunjukkan bahwa kebijakan penerimaan berkas tetap berlaku untuk sumber kedua melalui konfigurasi \textit{datasource}, bukan melalui logika khusus di dalam \textit{plugin}.
Dengan demikian, evaluasi ini memberikan bukti implementasi bahwa penambahan Google Drive dilakukan melalui \textit{plugin} dan konfigurasi sumber, sementara alur inti, penyimpanan, dan deduplikasi tetap digunakan.

\section{Evaluasi \textit{Legal Parser}}
\label{sec:evaluasi-legal-parser}

Evaluasi \textit{legal parser} dibagi menjadi tiga eksperimen yang saling berurutan.
Eksperimen pertama menilai \textit{classifier} yang menentukan apakah setiap halaman diproses melalui ekstraksi teks \textit{native} atau OCR.
Eksperimen kedua menguji toleransi penggabungan fragmen teks pada korpus yang seluruh halamannya diperkirakan bertipe \texttt{VECTOR}.
Eksperimen ketiga membandingkan model OCR dan penggunaan deteksi \textit{layout} pada korpus yang seluruh halamannya diperkirakan bertipe \texttt{SCAN}.

\subsection{Eksperimen Klasifikasi Halaman \textit{Vector} dan \textit{Scan}}
\label{sec:eksperimen-klasifikasi-halaman}

\subsubsection{Tujuan dan Pertanyaan Evaluasi}
\label{sec:tujuan-evaluasi-klasifikasi-halaman}

Eksperimen ini menetapkan aturan untuk mengarahkan setiap halaman PDF ke jalur ekstraksi yang sesuai.
Halaman yang memiliki teks \textit{native} yang memadai seharusnya tidak dikirim ke OCR karena OCR menambah biaya komputasi dan dapat memperkenalkan kesalahan pengenalan.
Sebaliknya, halaman hasil pindai yang dilewatkan ke ekstraksi \textit{native} dapat menghasilkan teks kosong atau teks yang tidak dapat digunakan untuk memulihkan struktur hukum.

Unit keputusan dibuat pada tingkat halaman karena satu dokumen dapat mencampurkan halaman \textit{native}, halaman pindai dengan \textit{text layer}, dan halaman pindai tanpa \textit{text layer}, sebagaimana dianalisis pada \autoref{sec:analisis-dokumen-peraturan}.
Pertanyaan evaluasi yang dijawab adalah sebagai berikut.

\begin{enumerate}
  \item Apakah ciri cakupan citra dan jumlah teks \textit{native} dapat membedakan halaman \texttt{VECTOR} dan \texttt{SCAN} dengan tingkat kesalahan yang rendah?
  \item Berapa nilai ambang cakupan citra dan jumlah teks minimum yang memberikan \textit{recall} tinggi pada kelas \texttt{SCAN}?
  \item Bagaimana penerapan ambang terpilih membentuk komposisi halaman dan dokumen pada seluruh korpus?
\end{enumerate}

\subsubsection{Rancangan Eksperimen}
\label{sec:rancangan-klasifikasi-halaman}

Korpus eksperimen adalah \textit{datasource} \path{bpk-peraturan-pendidikan} yang terdiri atas 8.745 dokumen aktif dengan 165.395 halaman.
Sampel dipilih secara acak seragam tanpa pengembalian pada tingkat halaman menggunakan seed \texttt{20260816}.
Sebanyak 500 halaman diberi label secara manual tanpa memperlihatkan prediksi atau ciri \textit{classifier} kepada pelabel.

Label \texttt{VECTOR} berarti teks \textit{native} pada halaman sudah cukup untuk digunakan, sedangkan label \texttt{SCAN} berarti halaman membutuhkan OCR untuk memperoleh isi yang dapat diandalkan.
Sebanyak 50 halaman ditampilkan kembali sehingga terdapat 550 presentasi untuk mengukur konsistensi pelabelan.
Pengulangan ini mengukur konsistensi terhadap halaman yang sama dan bukan kesepakatan antarpenilai.

\textit{Classifier} menggunakan dua parameter.
Parameter pertama adalah \path{scan_image_coverage_threshold}, yaitu rasio luas citra terbesar terhadap luas halaman.
Parameter kedua adalah \path{scan_min_text_threshold}, yaitu jumlah karakter teks \textit{native} setelah spasi di luar isi dibuang.
Seluruh 2.121 kombinasi pada grid dievaluasi, dengan cakupan citra 0--1 pada interval 0,01 dan teks minimum 0--200 karakter pada interval 10 karakter.

Metrik yang dihitung adalah akurasi, \textit{recall} dan \textit{precision} kelas \texttt{SCAN}, \textit{macro F1}, serta rasio halaman yang diarahkan ke OCR.
Kandidat diprioritaskan berdasarkan \textit{recall} \texttt{SCAN}, kemudian \textit{macro F1}, \textit{precision} \texttt{SCAN}, dan jumlah halaman yang diarahkan ke OCR.

\subsubsection{Hasil Evaluasi pada Sampel Berlabel}
\label{sec:hasil-klasifikasi-halaman}

Kandidat yang dipilih menggunakan ambang cakupan citra 0,50 dan ambang teks minimum 100 karakter.
Hasil evaluasinya dirangkum pada \autoref{tab:hasil-klasifikasi-halaman}.

\begin{table}[htb]
  \centering
  \footnotesize
  \caption{Hasil evaluasi \textit{classifier} halaman}
  \label{tab:hasil-klasifikasi-halaman}
  \begin{tabular}{@{}lr@{}}
    \toprule
    \textbf{Metrik} & \textbf{Nilai} \\
    \midrule
    Ambang cakupan citra & 0,50 \\
    Ambang teks minimum & 100 karakter \\
    Akurasi & 98,8\% \\
    \textit{SCAN recall} & 100,0\% \\
    \textit{SCAN precision} & 97,1\% \\
    \textit{Macro F1} & 98,8\% \\
    Rasio \textit{routing} OCR pada sampel & 42,0\% \\
    Konsistensi presentasi ulang & 98,0\% \\
    Mismatch terhadap implementasi & 0 \\
    \bottomrule
  \end{tabular}
\end{table}

Tidak ada halaman berlabel \texttt{SCAN} yang diprediksi sebagai \texttt{VECTOR}.
Sebanyak 204 halaman berlabel \texttt{SCAN} diprediksi sebagai \texttt{SCAN}, sedangkan enam halaman berlabel \texttt{VECTOR} ikut diarahkan ke OCR.
Hasil tersebut menunjukkan bahwa ambang yang dipilih lebih mengutamakan pencegahan halaman pindai lolos ke jalur \textit{native} daripada meminimalkan seluruh pemanggilan OCR.

\begin{table}[htb]
  \centering
  \footnotesize
  \caption{Matriks kebingungan \textit{classifier} pada sampel berlabel}
  \label{tab:confusion-klasifikasi-halaman}
  \begin{tabular}{@{}lrr@{}}
    \toprule
    & \textbf{Prediksi \texttt{SCAN}} & \textbf{Prediksi \texttt{VECTOR}} \\
    \midrule
    Label \texttt{SCAN} & 204 & 0 \\
    Label \texttt{VECTOR} & 6 & 290 \\
    \bottomrule
  \end{tabular}
\end{table}

\begin{figure}[htb]
  \centering
  \includegraphics[width=0.90\textwidth]{experiments/parser-classifier-threshold-grid.png}
  \caption{Evaluasi grid ambang \textit{classifier}; bintang menandai kandidat terpilih}
  \label{fig:classifier-threshold-grid}
\end{figure}

\subsubsection{Penerapan Ambang pada Seluruh Korpus}
\label{sec:komposisi-korpus-klasifikasi}

Ambang terpilih kemudian diterapkan pada seluruh halaman dalam \textit{datasource}.
Tahap ini merupakan analisis komposisi korpus berdasarkan prediksi \textit{classifier}, bukan evaluasi kedua dengan label manual.
Hasil \textit{routing} pada tingkat halaman terdiri atas 71.271 halaman \texttt{SCAN} atau 43,09\% dan 94.124 halaman \texttt{VECTOR} atau 56,91\%.

Pada tingkat dokumen, 3.858 dokumen hanya terdiri atas halaman \texttt{SCAN}, 4.459 dokumen hanya terdiri atas halaman \texttt{VECTOR}, dan 428 dokumen memiliki campuran kedua jenis halaman.
Kohort tersebut masing-masing mencakup 64.586, 79.017, dan 21.792 halaman.
Keberadaan 428 dokumen campuran menjadi alasan penting untuk mempertahankan keputusan \textit{routing} pada tingkat halaman, bukan pada tingkat dokumen.

\begin{table}[htb]
  \centering
  \footnotesize
  \caption{Komposisi korpus berdasarkan hasil \textit{routing} halaman}
  \label{tab:komposisi-korpus-parser}
  \begin{tabular}{@{}lrrr@{}}
    \toprule
    \textbf{Komposisi dokumen} & \textbf{Dokumen} & \textbf{Halaman} & \textbf{Proporsi halaman} \\
    \midrule
    \texttt{ALL\_SCAN} & 3.858 & 64.586 & 39,05\% \\
    \texttt{ALL\_VECTOR} & 4.459 & 79.017 & 47,77\% \\
    \texttt{MIXED} & 428 & 21.792 & 13,18\% \\
    \midrule
    \textbf{Total} & \textbf{8.745} & \textbf{165.395} & \textbf{100,00\%} \\
    \bottomrule
  \end{tabular}
\end{table}

\begin{figure}[htb]
  \centering
  \subfloat[Distribusi jenis halaman]{\includegraphics[width=0.47\textwidth]{experiments/parser-classifier-page-kind.png}}
  \hfill
  \subfloat[Komposisi dokumen]{\includegraphics[width=0.47\textwidth]{experiments/parser-classifier-document-composition.png}}
  \caption{Komposisi hasil \textit{routing} pada seluruh korpus peraturan}
  \label{fig:classifier-corpus-composition}
\end{figure}

\subsubsection{Kesimpulan Eksperimen}
\label{sec:kesimpulan-klasifikasi-halaman}

Eksperimen menunjukkan bahwa dua ciri struktural PDF dapat digunakan untuk membuat keputusan \textit{routing} pada tingkat halaman.
Ambang cakupan citra 0,50 dan teks minimum 100 karakter dipilih karena mempertahankan \textit{recall} kelas \texttt{SCAN} sebesar 100,0\% pada sampel berlabel dengan \textit{macro F1} sebesar 98,8\%.
Ambang tersebut digunakan sebagai kontrol tetap pada eksperimen \textit{legal parser} berikutnya.
Penerapan pada seluruh korpus memperkirakan bahwa 43,09\% halaman akan menggunakan jalur OCR dan 13,18\% halaman berada dalam dokumen campuran.

\subsection{Eksperimen \textit{Legal Parser} pada Korpus \texttt{ALL\_VECTOR}}
\label{sec:eksperimen-legal-parser-vector}

\subsubsection{Tujuan dan Pertanyaan Evaluasi}
\label{sec:tujuan-evaluasi-legal-parser-vector}

Eksperimen ini mengevaluasi jalur ekstraksi teks \textit{native} pada dokumen yang seluruh halamannya diprediksi sebagai \texttt{VECTOR}.
Jalur tersebut membaca fragmen teks beserta posisi geometrinya, menggabungkan fragmen menjadi baris, lalu meneruskan baris tersebut ke analisis penanda dan pemulihan struktur hukum.

Pertanyaan evaluasi yang dijawab adalah sebagai berikut.

\begin{enumerate}
  \item Bagaimana perubahan toleransi penggabungan vertikal dan horizontal memengaruhi keutuhan struktur hasil \textit{legal parser}?
  \item Nilai toleransi mana yang memberikan \textit{clean rate} terbaik pada korpus dokumen bertipe \texttt{VECTOR}?
  \item Apakah jalur \textit{native} dapat memproses seluruh korpus tanpa kegagalan terminal dan dengan biaya yang sesuai untuk pemrosesan korpus?
\end{enumerate}

\subsubsection{Rancangan Eksperimen}
\label{sec:rancangan-legal-parser-vector}

Eksperimen menggunakan sensus seluruh 4.459 dokumen \texttt{ALL\_VECTOR} yang mencakup 79.017 halaman karena jalur \textit{native} cukup murah secara komputasi dibandingkan jalur OCR sehingga tidak memerlukan pengambilan sampel.
Dokumen \texttt{MIXED} tidak disertakan karena halaman pindainya akan menggunakan jalur OCR.

Sebanyak 48 profil dijalankan melalui kombinasi dua parameter.
Parameter vertikal \path{merge_y_tolerance} diuji pada nilai 0, 2, 4, 6, 8, 10, 12, dan 16 pt.
Parameter horizontal \path{merge_x_tolerance} diuji pada nilai 0, 0,5, 1, 2, 4, dan 6 pt.
Ambang \textit{routing} hasil eksperimen sebelumnya tetap menggunakan cakupan citra 0,50 dan teks minimum 100 karakter.

Dokumen disebut bersih apabila berhasil menghasilkan dokumen kanonik tanpa peringatan.
Makna kode peringatan yang digunakan pada evaluasi telah ditetapkan pada \autoref{sec:solusi-verifikasi-ekstraksi}.
Klasifikasi kode peringatan pada eksperimen dirangkum pada \autoref{tab:klasifikasi-warning-legal-parser}.
Dokumen tanpa peringatan dikategorikan sebagai bersih.
Jika terdapat kedua jenis peringatan, kategori \textit{severe} diprioritaskan.
Kegagalan dihitung terpisah karena dokumen tidak menghasilkan keluaran \textit{parser}.
Definisi \textit{outcome} ini juga digunakan pada eksperimen jalur OCR pada subbab berikutnya.

\begin{table}[htb]
  \centering
  \scriptsize
  \caption{Klasifikasi kode peringatan pada eksperimen \textit{legal parser}}
  \label{tab:klasifikasi-warning-legal-parser}
  \begin{tabularx}{\textwidth}{@{}p{2.0cm}>{\raggedright\arraybackslash}p{6.4cm}>{\raggedright\arraybackslash}X@{}}
    \toprule
    \textbf{Kategori} & \textbf{Kode} & \textbf{Interpretasi} \\
    \midrule
    \textit{Advisory} &
    \path{sequence_gap}, \path{revise_pasal_count}, \path{BODY_STRUCTURAL_MARKER_IN_REVISION} &
    Dokumen kanonik tetap terbentuk, tetapi terdapat ketidakteraturan sumber atau konteks pembacaan yang perlu ditinjau. \\
    \textit{Severe} &
    \path{empty_leaf}, \path{revise_pasal_no_children}, \path{location_invalid}, \path{pembukaan_incomplete}, \path{empty_batang_tubuh} &
    Terdapat risiko lebih besar bahwa bagian struktur hilang, kosong, atau tidak terbentuk sesuai aturan domain. \\
    \bottomrule
  \end{tabularx}
\end{table}

\begin{table}[htb]
  \centering
  \footnotesize
  \caption{Rancangan eksperimen \textit{legal parser} pada korpus \texttt{ALL\_VECTOR}}
  \label{tab:rancangan-vector-parser}
  \begin{tabular}{@{}lr@{}}
    \toprule
    \textbf{Komponen} & \textbf{Nilai} \\
    \midrule
    Dokumen & 4.459 \\
    Halaman & 79.017 \\
    Jumlah profil & 48 \\
    Level \path{merge_y_tolerance} & 0, 2, 4, 6, 8, 10, 12, 16 pt \\
    Level \path{merge_x_tolerance} & 0, 0,5, 1, 2, 4, 6 pt \\
    Jalur ekstraksi & Teks \textit{native} \\
    \bottomrule
  \end{tabular}
\end{table}

\subsubsection{Hasil dan Analisis Pengaruh Toleransi Penggabungan}
\label{sec:hasil-legal-parser-vector}

Hasil seluruh profil memperlihatkan bahwa tidak ada dokumen yang gagal pada 48 profil.
Dengan demikian, terdapat 214.032 pemrosesan dokumen tanpa kegagalan terminal.
Perbedaan antarkonfigurasi terutama terlihat pada proporsi dokumen yang bersih dan jumlah peringatan struktural.

Profil \path{y6-x0.5} dipilih sebagai profil operasional.
Profil tersebut menghasilkan 3.216 dokumen bersih, 1.026 dokumen \textit{advisory}, 217 dokumen \textit{severe}, dan tidak memiliki kegagalan.
\textit{Clean rate}-nya adalah 72,1\%.
Perbandingan \textit{clean rate} pada seluruh 48 profil ditampilkan pada \autoref{fig:vector-clean-rate-grid}.

\begin{table}[htb]
  \centering
  \scriptsize
  \caption{Hasil profil operasional \textit{legal parser} pada korpus \texttt{ALL\_VECTOR}}
  \label{tab:hasil-vector-parser}
  \begin{tabular}{@{}lrrrrrr@{}}
    \toprule
    \textbf{Profil} & \textbf{Bersih} & \textbf{\textit{Advisory}} & \textbf{\textit{Severe}} & \textbf{Gagal} & \textbf{\textit{Clean rate}} & \textbf{Peringatan} \\
    \midrule
    \path{y6-x0.5} & 3.216 & 1.026 & 217 & 0 & 72,1\% & 3.060 \\
    \bottomrule
  \end{tabular}
\end{table}

\begin{figure}[htb]
  \centering
  \includegraphics[width=0.88\textwidth]{experiments/parser-vector-clean-rate-grid.png}
  \caption{\textit{Clean rate} pada grid toleransi penggabungan jalur \textit{native}}
  \label{fig:vector-clean-rate-grid}
\end{figure}

Rata-rata marginal adalah rata-rata \textit{clean rate} pada seluruh nilai parameter lain.
Dengan demikian, rata-rata untuk satu nilai \path{merge_y_tolerance} dihitung dari enam nilai \path{merge_x_tolerance}, sedangkan rata-rata untuk satu nilai \path{merge_x_tolerance} dihitung dari delapan nilai \path{merge_y_tolerance}.
Angka marginal digunakan untuk membaca kecenderungan satu parameter secara terpisah dan bukan untuk menggantikan hasil pada satu sel konfigurasi tertentu.
Rata-rata marginal \path{merge_y_tolerance} meningkat dari 63,3\% pada 0 pt menjadi 70,9\% pada 6 pt.
Nilainya berada pada sekitar 70,9\% sampai 10 pt, lalu menurun menjadi 62,5\% pada 16 pt.
Pada parameter horizontal, nilai 0,5 pt menghasilkan rata-rata marginal tertinggi sebesar 69,9\%, sedangkan nilai 4 dan 6 pt menghasilkan 66,9\%.
Artinya, toleransi vertikal antara 4--10 pt relatif stabil pada korpus ini, sementara toleransi horizontal yang terlalu besar cenderung menurunkan \textit{clean rate}.

\begin{table}[htb]
  \centering
  \scriptsize
  \caption{\textit{Clean rate} marginal berdasarkan toleransi penggabungan}
  \label{tab:efek-toleransi-vector}
  \begin{tabular}{@{}llr@{}}
    \toprule
  \textbf{Parameter} & \textbf{Nilai} & \textbf{\textit{Clean rate} marginal} \\
    \midrule
    \path{merge_y_tolerance} & 0 pt & 63,3\% \\
    & 2 pt & 70,5\% \\
    & 4 pt & 70,8\% \\
    & 6 pt & 70,9\% \\
    & 8 pt & 70,9\% \\
    & 10 pt & 70,9\% \\
    & 12 pt & 70,7\% \\
    & 16 pt & 62,5\% \\
    \midrule
    \path{merge_x_tolerance} & 0 pt & 69,8\% \\
    & 0,5 pt & 69,9\% \\
    & 1 pt & 69,8\% \\
    & 2 pt & 69,6\% \\
    & 4 pt & 66,9\% \\
    & 6 pt & 66,9\% \\
    \bottomrule
  \end{tabular}
\end{table}

Komposisi \textit{outcome} untuk seluruh 48 profil ditampilkan pada \autoref{fig:vector-outcome-composition}.
Gambar tersebut memperlihatkan bahwa perubahan toleransi mengubah proporsi dokumen bersih, \textit{advisory}, dan \textit{severe}, meskipun seluruh profil tetap menghasilkan keluaran \textit{parser} tanpa kegagalan.

\begin{figure}[H]
  \centering
  \includegraphics[width=0.72\textwidth]{experiments/parser-vector-outcome-composition.png}
  \caption{Komposisi hasil pemrosesan pada grid profil \textit{legal parser} jalur \textit{native}}
  \label{fig:vector-outcome-composition}
\end{figure}

Pada profil \path{y6-x0.5}, peringatan yang paling banyak adalah \path{sequence_gap} sebanyak 2.171 kejadian.
Peringatan tersebut tidak selalu menunjukkan kesalahan \textit{parser} karena dokumen sumber sering kali memang melewati nomor tertentu.
Peringatan lain yang muncul adalah \path{BODY_STRUCTURAL_MARKER_IN_REVISION} sebanyak 494 kejadian, \path{pembukaan_incomplete} sebanyak 213 kejadian, dan \path{empty_leaf} sebanyak 131 kejadian.
Dua kode lain yang muncul adalah \path{revise_pasal_no_children} sebanyak 49 kejadian dan \path{revise_pasal_count} sebanyak 2 kejadian.
Makna keenam kode tersebut mengikuti definisi pada \autoref{tab:kode-peringatan-legal-parser}.
Distribusi seluruh kode peringatan pada profil-profil yang diuji ditampilkan pada \autoref{fig:vector-warning-codes}.

\begin{figure}[H]
  \centering
  \includegraphics[width=0.92\textwidth]{experiments/parser-vector-warning-codes.png}
  \caption{Peringatan \textit{parser} pada jalur \textit{native}}
  \label{fig:vector-warning-codes}
\end{figure}

Pengukuran \textit{throughput} dilakukan pada profil \path{y6-x0.5} dalam eksekusi \textit{standalone} sehingga tidak bercampur dengan kontensi dari profil lain.
Profil tersebut memproses 552 dokumen per menit atau 163 halaman per detik.
Median waktu pemrosesan satu dokumen adalah 72 ms, sedangkan tambahan waktu per halaman setelah halaman pertama adalah sekitar 5,0 ms.
Perkembangan pemrosesan dan hubungan waktu dengan jumlah halaman ditampilkan pada \autoref{fig:vector-throughput}.

\begin{figure}[H]
  \centering
  \includegraphics[width=0.92\textwidth]{experiments/parser-vector-throughput.png}
  \caption{\textit{Throughput} profil \texttt{y6-x0.5}}
  \label{fig:vector-throughput}
\end{figure}

\subsubsection{Analisis dan Kesimpulan Eksperimen}
\label{sec:kesimpulan-legal-parser-vector}

Toleransi vertikal yang terlalu ketat dapat memecah fragmen teks yang seharusnya berada pada satu baris, sedangkan toleransi yang terlalu longgar dapat menggabungkan fragmen dari baris yang berbeda.
Pola tersebut terlihat pada penurunan \textit{clean rate} di 0 pt dan 16 pt.
Toleransi horizontal yang besar juga menurunkan \textit{clean rate} karena lebih banyak fragmen yang dapat digabungkan tanpa pemisah yang semestinya.
Profil \path{y6-x0.5} berada pada \textit{plateau clean rate} terbaik dan dipilih sebagai konfigurasi \textit{native} operasional.
\textit{Clean rate} dan peringatan pada eksperimen ini merupakan indikator operasional berbasis pemeriksaan \textit{parser}, bukan pengukuran terhadap anotasi struktur hukum yang dibuat secara manual.
Seluruh profil berhasil memproses korpus tanpa kegagalan, sedangkan profil operasional mencapai 552 dokumen per menit.

\subsection{Eksperimen \textit{Legal Parser} pada Korpus \texttt{ALL\_SCAN}}
\label{sec:eksperimen-legal-parser-ocr}

\subsubsection{Tujuan dan Pertanyaan Evaluasi}
\label{sec:tujuan-evaluasi-legal-parser-ocr}

Eksperimen ini mengevaluasi jalur OCR pada dokumen hukum hasil pindai.
Istilah \textit{jalur OCR} merujuk pada tahap pengambilan teks halaman sebelum baris tersebut diproses oleh lexer, segmenter, dan pemeriksa struktur \textit{legal parser}.
Dengan demikian, eksperimen tidak hanya mengukur layanan OCR secara terpisah, tetapi mengukur dampaknya terhadap keluaran struktur hukum.

Pertanyaan evaluasi yang dijawab adalah sebagai berikut.

\begin{enumerate}
  \item Model OCR mana yang menghasilkan proporsi dokumen bersih paling tinggi pada korpus peraturan hasil pindai?
  \item Apakah deteksi \textit{layout} meningkatkan keutuhan struktur atau justru menambah gangguan pada urutan baca?
  \item Bagaimana perubahan model dan deteksi \textit{layout} memengaruhi waktu pemrosesan, \textit{throughput}, dan kegagalan pengenalan?
\end{enumerate}

\subsubsection{Rancangan Eksperimen}
\label{sec:rancangan-legal-parser-ocr}

Sampel dipilih dari 2.577 dokumen \texttt{ALL\_SCAN} yang memiliki paling banyak 15 halaman.
Sebanyak 250 dokumen dipilih secara acak tanpa pengembalian menggunakan seed \texttt{20260816}.
Sampel tersebut mencakup 2.249 halaman dan digunakan kembali pada keempat profil agar perbandingan dilakukan pada dokumen dan halaman yang sama.

Eksperimen menggunakan rancangan faktorial 2 x 2 yang mengombinasikan dua model OCR dan dua pilihan deteksi \textit{layout}.
Model OCR yang diuji adalah GLM-OCR dan PaddleOCR-VL 1.6.
Deteksi \textit{layout} menggunakan PP-DocLayoutV3 atau dinonaktifkan.

Parameter faktor dan kontrol eksperimen dirangkum pada \autoref{tab:kontrol-ocr-parser}.

Metrik yang dicatat adalah jumlah dan proporsi dokumen bersih, \textit{advisory}, \textit{severe}, dan gagal, jumlah peringatan berdasarkan kode, jenis kegagalan, waktu per halaman, \textit{throughput} dokumen, \textit{throughput} halaman, serta jumlah permintaan OCR per halaman.

\begin{table}[htb]
  \centering
  \scriptsize
  \caption{Parameter eksperimen \textit{legal parser} jalur OCR}
  \label{tab:kontrol-ocr-parser}
  \begin{tabularx}{\textwidth}{@{}p{4.2cm}>{\raggedright\arraybackslash}p{3.6cm}>{\raggedright\arraybackslash}X@{}}
    \toprule
    \textbf{Parameter} & \textbf{Nilai} & \textbf{Peran} \\
    \midrule
    Model OCR & GLM-OCR; PaddleOCR-VL 1.6 & Faktor eksperimen \\
    Prompt OCR & \texttt{Text Recognition:}; \texttt{OCR:} & Prompt sesuai model OCR \\
    Model \textit{layout} & \path{none}; PP-DocLayoutV3 & Faktor eksperimen \\
    \textit{Region} dikecualikan & \path{chart}, \path{image}, \path{seal} & Diterapkan saat PP-DocLayoutV3 digunakan \\
    Resolusi rasterisasi & 150 DPI & Kontrol kualitas citra halaman \\
    \path{scan_image_coverage_threshold} & 0,5 & Ambang \textit{routing} dari eksperimen classifier \\
    \path{scan_min_text_threshold} & 100 & Ambang teks \textit{native} dari eksperimen classifier \\
    \path{max_tokens} & 4.096 & Batas keluaran satu permintaan OCR \\
    \path{temperature} & 0 & Decoding deterministik \\
    \path{repetition_penalty} & 1,1 & Mengurangi pengulangan keluaran \\
    \path{top_p} & 0,00001 & Membatasi kandidat probabilitas token \\
    \path{top_k} & 1 & Memilih token dengan probabilitas tertinggi \\
    \path{max_concurrent_requests} & 4 & Batas permintaan konkuren \\
    Ambang skor \textit{layout} & 0,3 & Kontrol deteksi region \\
    \bottomrule
  \end{tabularx}
\end{table}

Parameter \path{merge_y_tolerance} dan \path{merge_x_tolerance} tidak menjadi faktor pada eksperimen ini karena jalur OCR membentuk baris dari keluaran model, bukan dari penggabungan fragmen teks berdasarkan geometri PDF.

\begin{table}[htb]
  \centering
  \scriptsize
  \caption{Profil eksperimen \textit{legal parser} jalur OCR}
  \label{tab:konfigurasi-ocr-parser}
  \begin{tabularx}{\textwidth}{@{}>{\raggedright\arraybackslash}X>{\raggedright\arraybackslash}X>{\raggedright\arraybackslash}X>{\raggedright\arraybackslash}X@{}}
    \toprule
    \textbf{Profil} & \textbf{Model OCR} & \textbf{\textit{Layout}} & \textbf{\textit{Region} yang dikecualikan} \\
    \midrule
    \path{glm-ocr} & GLM-OCR & Tidak ada & Tidak ada \\
    \shortstack[l]{\texttt{glm-ocr-}\\\texttt{pp-doclayout-v3}} & GLM-OCR & PP-DocLayoutV3 & \path{chart}, \path{image}, \path{seal} \\
    \path{paddle-ocr} & PaddleOCR-VL 1.6 & Tidak ada & Tidak ada \\
    \shortstack[l]{\texttt{paddle-ocr-}\\\texttt{pp-doclayout-v3}} & PaddleOCR-VL 1.6 & PP-DocLayoutV3 & \path{chart}, \path{image}, \path{seal} \\
    \bottomrule
  \end{tabularx}
\end{table}

\subsubsection{Hasil Keutuhan Struktur}
\label{sec:hasil-legal-parser-ocr}

Profil \path{glm-ocr} tanpa deteksi \textit{layout} menghasilkan proporsi dokumen bersih tertinggi, yaitu 173 dari 250 dokumen atau 69,2\%.
Profil tersebut juga menghasilkan sembilan dokumen \textit{severe} dan satu kegagalan.
Hasil keempat profil dirangkum pada \autoref{tab:hasil-ocr-parser}.

\begin{table}[htb]
  \centering
  \scriptsize
  \caption{Hasil keutuhan struktur pada eksperimen \textit{legal parser} jalur OCR}
  \label{tab:hasil-ocr-parser}
  \begin{tabular}{@{}lrrrrrr@{}}
    \toprule
    \textbf{Profil} & \textbf{Bersih} & \textbf{\textit{Advisory}} & \textbf{\textit{Severe}} & \textbf{Gagal} & \textbf{\textit{Clean rate}} & \textbf{\textit{Failure rate}} \\
    \midrule
    \path{glm-ocr} & 173 & 67 & 9 & 1 & 69,2\% & 0,4\% \\
    \path{glm-ocr-pp-doclayout-v3} & 37 & 144 & 67 & 2 & 14,8\% & 0,8\% \\
    \path{paddle-ocr} & 97 & 122 & 20 & 11 & 38,8\% & 4,4\% \\
    \path{paddle-ocr-pp-doclayout-v3} & 68 & 109 & 67 & 6 & 27,2\% & 2,4\% \\
    \bottomrule
  \end{tabular}
\end{table}

\begin{figure}[htb]
  \centering
  \includegraphics[width=0.88\textwidth]{experiments/parser-ocr-outcome-composition.png}
  \caption{Komposisi hasil pemrosesan pada eksperimen \textit{legal parser} jalur OCR}
  \label{fig:ocr-outcome-composition}
\end{figure}

Penambahan PP-DocLayoutV3 menurunkan \textit{clean rate} pada kedua model OCR.
Pada GLM-OCR, \textit{clean rate} turun dari 69,2\% menjadi 14,8\%, sedangkan pada PaddleOCR-VL 1.6 turun dari 38,8\% menjadi 27,2\%.
Peringatan \path{sequence_gap} juga meningkat pada GLM-OCR dari 91 menjadi 598 dan pada PaddleOCR dari 269 menjadi 364 ketika deteksi \textit{layout} diaktifkan.
Pola tersebut menunjukkan bahwa peningkatan pemrosesan berbasis region tidak otomatis menghasilkan urutan baca yang lebih sesuai untuk struktur peraturan.

\begin{figure}[htb]
  \centering
  \includegraphics[width=0.92\textwidth]{experiments/parser-ocr-warning-codes.png}
  \caption{Peringatan \textit{parser} pada jalur OCR}
  \label{fig:ocr-warning-codes}
\end{figure}

\subsubsection{Biaya Pemrosesan dan Kegagalan OCR}
\label{sec:biaya-kegagalan-legal-parser-ocr}

PaddleOCR tanpa deteksi \textit{layout} memiliki waktu median terendah, yaitu 3,87 detik per halaman, dengan \textit{throughput} 1,69 dokumen per menit dan 0,240 halaman per detik.
GLM-OCR tanpa deteksi \textit{layout} membutuhkan 9,54 detik per halaman dan menghasilkan 0,68 dokumen per menit.
Namun, deteksi \textit{layout} mempercepat GLM-OCR menjadi 5,49 detik per halaman, sementara pada PaddleOCR justru memperlambatnya menjadi 6,32 detik per halaman.

Deteksi \textit{layout} mengganti satu permintaan OCR halaman penuh menjadi sekitar 17 permintaan \textit{region} per halaman.
Karena itu, dampaknya terhadap waktu bergantung pada model OCR dan biaya penjadwalan region.

\begin{table}[htb]
  \centering
  \scriptsize
  \caption{Biaya pemrosesan eksperimen \textit{legal parser} jalur OCR}
  \label{tab:biaya-ocr-parser}
  \begin{tabularx}{\textwidth}{@{}>{\raggedright\arraybackslash}Xcccc@{}}
    \toprule
    \textbf{Profil} & \shortstack{\textbf{Median}\\\textbf{detik/halaman}} & \shortstack{\textbf{Dokumen/}\\\textbf{menit}} & \shortstack{\textbf{Halaman/}\\\textbf{detik}} & \shortstack{\textbf{Permintaan OCR/}\\\textbf{halaman}} \\
    \midrule
    \path{glm-ocr} & 9,54 & 0,68 & 0,101 & 1,00 \\
    \shortstack[l]{\texttt{glm-ocr-}\\\texttt{pp-doclayout-v3}} & 5,49 & 1,17 & 0,174 & 17,15 \\
    \path{paddle-ocr} & 3,87 & 1,69 & 0,240 & 1,00 \\
    \shortstack[l]{\texttt{paddle-ocr-}\\\texttt{pp-doclayout-v3}} & 6,32 & 1,02 & 0,149 & 17,20 \\
    \bottomrule
  \end{tabularx}
\end{table}

Seluruh 20 kegagalan pada keempat profil terjadi pada tahap \path{ocr.recognition} ketika keluaran mencapai batas 4.096 token.
Jumlah kegagalan berturut-turut adalah satu pada GLM-OCR, dua pada GLM-OCR dengan \textit{layout}, sebelas pada PaddleOCR, dan enam pada PaddleOCR dengan \textit{layout}.
Kegagalan tersebut tidak dianggap dapat diselesaikan dengan pengulangan biasa karena decoding yang deterministik akan mencapai batas yang sama.

\begin{figure}[htb]
  \centering
  \subfloat[Durasi dan \textit{throughput}]{\includegraphics[width=0.47\textwidth]{experiments/parser-ocr-duration-throughput.png}}
  \hfill
  \subfloat[Kegagalan pengenalan]{\includegraphics[width=0.47\textwidth]{experiments/parser-ocr-failures.png}}
  \caption{Biaya dan kegagalan pada eksperimen \textit{legal parser} jalur OCR}
  \label{fig:ocr-cost-failures}
\end{figure}

\subsubsection{Analisis dan Kesimpulan Eksperimen}
\label{sec:kesimpulan-legal-parser-ocr}

Berdasarkan keutuhan struktur, \path{glm-ocr} tanpa \textit{layout} dipilih sebagai konfigurasi operasional karena menghasilkan \textit{clean rate} tertinggi dan jumlah kegagalan paling rendah.
PP-DocLayoutV3 menurunkan \textit{clean rate} pada kedua model, meskipun mempercepat GLM-OCR dari 9,54 menjadi 5,49 detik per halaman dan memperlambat PaddleOCR dari 3,87 menjadi 6,32 detik per halaman.
PaddleOCR tanpa \textit{layout} merupakan konfigurasi tercepat, tetapi mutu struktur yang dihasilkan lebih rendah.
Dengan demikian, eksperimen memilih GLM-OCR tanpa \textit{layout} untuk prioritas keutuhan struktur dan mencatat batas 4.096 token sebagai penyebab seluruh kegagalan OCR.

\section{Evaluasi \textit{Generic Outline Parser} pada Domain \textit{User Manual}}
\label{sec:evaluasi-generic-outline-customer-service}

Evaluasi ini menguji \textit{generic outline parser} pada manual produk NETGEAR yang digunakan sebagai sumber pengetahuan domain \textit{User Manual}.
Berbeda dari eksperimen \textit{legal parser}, evaluasi ini merupakan \textit{coverage check} dengan satu konfigurasi default dan tidak membandingkan beberapa model atau nilai parameter.

\subsection{Tujuan dan Pertanyaan Evaluasi}
\label{sec:tujuan-evaluasi-generic-outline}

Evaluasi bertujuan menilai apakah \textit{bookmark outline} pada manual produk dapat dipetakan menjadi hierarki dokumen kanonik.
Evaluasi juga memeriksa keberhasilan penautan setiap \textit{bookmark} ke baris judul dan kemampuan parser menghasilkan keluaran tanpa peringatan atau kegagalan.

Pertanyaan evaluasi yang dijawab adalah sebagai berikut.

\begin{enumerate}
  \item Apakah seluruh manual NETGEAR pada \textit{datasource} dapat diproses oleh \textit{generic outline parser}?
  \item Apakah setiap \textit{bookmark} dapat ditautkan ke judul bagian yang sesuai?
  \item Apakah struktur \textit{bookmark} dapat dimaterialisasi menjadi hierarki blok kanonik?
  \item Apakah bagian yang dikecualikan seperti \textit{Contents} tidak mengganggu pembentukan struktur?
\end{enumerate}

\subsection{Rancangan Evaluasi}
\label{sec:rancangan-evaluasi-generic-outline}

Evaluasi menggunakan seluruh 9 dokumen aktif pada \textit{datasource} \texttt{netgear\_products}, yang mencakup 1.400 halaman dan 1.566 \textit{bookmark}.
Dokumen diurutkan berdasarkan nama berkas sehingga pengulangan evaluasi menggunakan himpunan yang sama tanpa seed dan tanpa pengambilan sampel.
Seluruh dokumen diproses dengan satu konfigurasi default.

\begin{table}[H]
  \centering
  \scriptsize
  \caption{Konfigurasi evaluasi \textit{generic outline parser}}
  \label{tab:konfigurasi-generic-outline}
  \begin{tabularx}{\textwidth}{@{}>{\raggedright\arraybackslash}p{5.1cm}>{\raggedright\arraybackslash}p{3.4cm}>{\raggedright\arraybackslash}X@{}}
    \toprule
    \textbf{Parameter} & \textbf{Nilai} & \textbf{Peran} \\
    \midrule
    \path{header_margin_ratio} & 0,075 & Proporsi tinggi halaman yang diperlakukan sebagai \textit{header} \\
    \path{footer_margin_ratio} & 0,90 & Batas awal area \textit{footer} \\
    \path{anchor_search_window_points} & 150 pt & Jendela pencarian judul di sekitar tujuan \textit{bookmark} \\
    \path{line_merge_tolerance_points} & 4 pt & Toleransi penggabungan fragmen teks menjadi satu baris \\
    \path{text_decoration} & \path{true} & Mempertahankan dekorasi teks sebagai \textit{markdown} \\
    \path{excluded_section_titles} & \path{contents}, \path{table of contents} & Bagian yang dibuang beserta turunannya dan dicatat sebagai \textit{metadata} \\
    \bottomrule
  \end{tabularx}
\end{table}

\textit{Bookmark} menjadi sumber hierarki karena manual produk tidak memiliki aturan penomoran domain yang dapat digunakan untuk menyimpulkan tingkat judul.
Judul \textit{bookmark} dicari pada jendela baris di sekitar tujuan \textit{bookmark}.
Jika pencocokan judul tidak ditemukan, parser menggunakan baris pertama pada atau setelah tujuan sebagai \textit{fallback} dan mencatat peringatan sesuai kode yang berlaku.
PDF tanpa \textit{outline} atau tanpa teks yang dapat diekstraksi ditolak sebagai kegagalan sumber.
Parser tidak menggunakan OCR atau deteksi \textit{layout}, dan tabel tidak dimodelkan sebagai struktur khusus.

\subsection{Hasil Pemrosesan}
\label{sec:hasil-pemrosesan-generic-outline}

Seluruh 9 manual berhasil diproses sebagai dokumen bersih.
Tidak terdapat dokumen dengan peringatan, tidak ada dokumen gagal, dan tidak ada dokumen yang ditolak karena \textit{outline} atau teks tidak tersedia.
Ringkasan hasil pemrosesan ditampilkan pada \autoref{tab:ringkasan-generic-outline}.

\begin{table}[H]
  \centering
  \footnotesize
  \caption{Ringkasan hasil evaluasi \textit{generic outline parser}}
  \label{tab:ringkasan-generic-outline}
  \begin{tabular}{@{}lr@{}}
    \toprule
    \textbf{Metrik} & \textbf{Nilai} \\
    \midrule
    Dokumen diproses & 9 \\
    Dokumen bersih & 9 \\
    Dokumen dengan peringatan & 0 \\
    Dokumen gagal & 0 \\
    Halaman & 1.400 \\
    \textit{Bookmark} & 1.566 \\
    \textit{Heading} terbentuk & 1.557 \\
    \textit{Anchor title-match} & 1.566/1.566 (100\%) \\
    Durasi total & 7,6 detik \\
    \bottomrule
  \end{tabular}
\end{table}

Sebanyak 1.566 \textit{bookmark} menghasilkan 1.557 \textit{heading}.
Selisih sembilan \textit{heading} berasal dari bagian \textit{Contents} yang dikecualikan pada masing-masing manual.
Bagian yang dikecualikan disimpan dalam \textit{metadata} \path{excluded_sections} dan tidak dihitung sebagai peringatan.

Perincian setiap manual ditampilkan pada \autoref{tab:hasil-perdokumen-generic-outline}.

\begin{table}[H]
  \centering
  \scriptsize
  \caption{Hasil per dokumen pada evaluasi \textit{generic outline parser}}
  \label{tab:hasil-perdokumen-generic-outline}
  \begin{tabularx}{\textwidth}{@{}>{\raggedright\arraybackslash}Xrrrrrr@{}}
    \toprule
    \textbf{Dokumen} & \textbf{Halaman} & \textbf{\textit{Bookmark}} & \textbf{\textit{Heading}} & \textbf{\textit{Text block}} & \textbf{\textit{Title-match}} & \textbf{Detik} \\
    \midrule
    \path{netgear-cm2000.pdf} & 26 & 29 & 28 & 29 & 100\% & 0,15 \\
    \path{netgear-lm1200.pdf} & 105 & 105 & 104 & 105 & 100\% & 0,50 \\
    \path{netgear-r6350.pdf} & 204 & 226 & 225 & 226 & 100\% & 1,08 \\
    \path{netgear-r7000.pdf} & 186 & 213 & 212 & 213 & 100\% & 0,97 \\
    \path{netgear-rax120.pdf} & 175 & 202 & 201 & 202 & 100\% & 1,08 \\
    \path{netgear-rax50.pdf} & 160 & 185 & 184 & 185 & 100\% & 0,87 \\
    \path{netgear-raxe500.pdf} & 169 & 192 & 191 & 192 & 100\% & 0,93 \\
    \path{netgear-rbk852.pdf} & 161 & 183 & 182 & 183 & 100\% & 0,85 \\
    \path{netgear-xr500.pdf} & 214 & 231 & 230 & 231 & 100\% & 1,17 \\
    \midrule
    \textbf{Total} & \textbf{1.400} & \textbf{1.566} & \textbf{1.557} & \textbf{1.566} & \textbf{100\%} & \textbf{7,60} \\
    \bottomrule
  \end{tabularx}
\end{table}

Jumlah \textit{text block} melebihi jumlah \textit{heading} sebanyak satu pada setiap dokumen.
Blok tambahan tersebut merupakan materi awal dokumen sebelum \textit{bookmark} pertama dan disimpan sebagai \textit{front matter}.

\subsection{Analisis Anchoring, Kinerja, dan Kesimpulan}
\label{sec:analisis-kesimpulan-generic-outline}

Seluruh 1.566 \textit{bookmark} berhasil dicocokkan dengan judulnya.
Hasil ini mendukung strategi pencarian judul pada jendela di sekitar tujuan \textit{bookmark}, bukan sekadar mengambil baris pertama setelah tujuan.
Pada konfigurasi ini, seluruh dokumen juga lolos pemeriksaan struktur tanpa \textit{empty node} atau rentang halaman yang tidak valid.

Evaluasi selesai dalam 7,6 detik, dengan \textit{throughput} 71 dokumen per menit dan 184 halaman per detik.
Median waktu pemrosesan adalah 0,93 detik per dokumen, sedangkan biaya tambahan berdasarkan jumlah halaman adalah sekitar 5,4 ms per halaman.

Dengan demikian, \textit{generic outline parser} berhasil memulihkan hierarki manual produk menjadi dokumen kanonik pada seluruh 9 dokumen yang diuji.
Anchoring \textit{bookmark} ke judul berhasil pada seluruh 1.566 \textit{bookmark}.
Pengecualian bagian \textit{Contents} bekerja sesuai konfigurasi dan tidak diperlakukan sebagai cacat struktur.
Hasil ini mendukung penggunaan kontrak \textit{parser plugin} yang sama untuk domain \textit{User Manual} dengan strategi ekstraksi berbasis \textit{outline}.
Evaluasi ini tidak digeneralisasikan ke seluruh PDF ber-\textit{bookmark} karena korpus hanya berasal dari satu produsen dan satu \textit{toolchain}.
